% ch-a4-09.tex —— 第 4 篇（Emerton–Gee–Hellmann）续片，底本 PDF 276–287 = 印刷页 263–274

\ysrc{263}
……映到 $G$ 的局部解析表示。我们回顾其构造的一部分。事实上，可以设想 Orlik--Strauch 的构造可推广为函子
\[
\mathcal F_P^G:\Rep_L^{\mathrm{an}}M\lra\Rep_L^{\mathrm{an}}G,\qquad
\mathcal F_P^G:D(\mathrm{an}\,M)\lra D(\mathrm{an}\,G),
\]
其行为与抛物诱导类似，见下文期望 6.2.22。

下面仍聚焦于 $P=B$ 的情形。以下用 $U(\mathfrak g_L)$ 表示 $\mathfrak g_L=(\operatorname{Lie}G)\otimes_{\Q_p}L$ 的泛包络代数（注意此处把 $G$ 看作 $\Q_p$-解析群，故 $\mathfrak g_L$ 同构于乘积 $\prod_{\tau:F\hookrightarrow L}\mathfrak{gl}_d$）。类似地，用 $\mathfrak t_L$ 与 $\mathfrak b_L$ 表示 $\Q_p$-解析群 $T$、$B$ 在 $L$ 上的李代数。以 $\mathcal O(U(\mathfrak g_L))$ 表示 BGG 范畴 $\mathcal O$：其对象是有限生成 $U(\mathfrak g_L)$-模，Cartan $\mathfrak t_L$ 在其上半单地作用，且作为 $U(\mathfrak b_L)$-模局部有限。以 $\mathcal O(U(\mathfrak g_L))^{\mathrm{alg}}$ 表示由 $\mathfrak t_L$ 的作用为代数的对象 $M$ 构成的子范畴。在该子范畴上可以把 $\mathfrak b_L$-作用积分，得到 $M$ 上的 $B$-作用。于是 $\mathcal O(U(\mathfrak g_L))^{\mathrm{alg}}$ 的每个对象 $M$ 都可看作由 $U(\mathfrak g_L)$ 与 $\mathcal D(B)$ 生成的子 $L$-代数 $\mathcal D(\mathfrak g_L,B)\subseteq\mathcal D(G)$ 上的模。由此得到正合反变函子
\[
F_B^G(-):\mathcal O(U(\mathfrak g_L))^{\mathrm{alg}}\lra\Rep^{\mathrm{an}}G,
\]
由 $V\mapsto(\mathcal D(G)\otimes_{\mathcal D(\mathfrak g,B)}V)'$ 给出；参看 [OS15, 3.4]。

\begin{rem}{6.2.18}\label{rem:a4-6-2-18}
在对偶侧，用 $\mathcal D(G)$-模来说，局部解析抛物诱导由
\[
\bigl((\Ind_B^G\delta)^{\mathrm{an}}\bigr)'_{\mathrm b}
=\mathcal D(G)\otimes_{\mathcal D(B)}(L_\delta)'_{\mathrm b}
\]
给出，其中 $L_\delta$ 表示 $T$ 通过特征标 $\delta$ 在 $L$ 上的 1 维局部解析表示，见 [OS15, Lemma 2.3]。现设 $\xi$ 是 $T$ 的代数特征标。Verma 模 $M(\xi)=U(\mathfrak g_L)\otimes_{U(\mathfrak b_L)}L_\xi$ 是 $\mathcal O(U(\mathfrak g_L))^{\mathrm{alg}}$ 的对象，并且
\[
\mathcal D(G)\otimes_{\mathcal D(\mathfrak g_L,B)}M(\xi)=\mathcal D(G)\otimes_{\mathcal D(B)}L_\xi.
\]
特别地，对代数特征标 $\xi$ 有
\[
F_B^G(M(-\xi))=\Ind_B^G(z^\xi)^{\mathrm{an}}.
\]
事实上，Orlik 与 Strauch [OS15, 4.4] 证明：函子 $F_B^G(-)$ 在某种意义上可以用 $T$ 的光滑特征标来扭。具体地说，对 $T$ 的任一光滑特征标 $\delta_{\mathrm{sm}}$，存在正合反变函子
\[
F_B^G(-,\delta_{\mathrm{sm}}):\mathcal O(U(\mathfrak g_L))^{\mathrm{alg}}\lra\Rep^{\mathrm{an}}G,
\]
使得对代数特征标 $\xi$，表示 $F_B^G(M(-\xi),\delta_{\mathrm{sm}})$ 就是局部解析诱导表示 $(\Ind_B^Gz^\xi\delta_{\mathrm{sm}})^{\mathrm{la}}$。
\end{rem}

我们指出它与光滑部分为抛物诱导的局部代数表示之间的关系。固定一个整的、支配的、正则的 Hodge--Tate 权 $\underline\lambda$，并用 $\xi$ 表示 $\mathfrak t_L$ 的相应特征标（通过定义 2.3.1），用 $\mathcal O(U(\mathfrak g_L))_{(\xi)}\subseteq\mathcal O(U(\mathfrak g_L))$ 表示范畴 $\mathcal O$ 中含 Verma 模 $M(\xi)$ 的区段（block）。则 $\mathcal O(U(\mathfrak g_L))_{(\xi)}\subseteq\mathcal O(U(\mathfrak g_L))^{\mathrm{alg}}$\tn{底本此处两处均作 $U(\mathfrak g)$（无下标 $L$），与同段其余记号不一致，译文统一为 $U(\mathfrak g_L)$。}。由于 $\lambda$ 正则，特征标 $\xi$ 自动是点正则的（dot-regular），即 $\xi$ 在点作用 $w\cdot\xi=w(\xi+\rho)-\rho$ 下的 Weyl 群轨道 $W\cdot\xi$ 恰含 $|W|$ 个元素。

\ysrc{264}
$\mathcal O(U(\mathfrak g_L))_{(\xi)}$ 的单对象是 Verma 模 $M(w\cdot\xi)$ 的单商 $L(w\cdot\xi)$。此外，$L(\xi)=W_{\underline\lambda}[1/p]$ 是定义 2.3.2 中定义的 $G$ 的不可约代数表示。

\begin{lem}{6.2.19}\label{lem:a4-6-2-19}
存在典范同构
\[
F_B^G(L(-\xi),\delta_{\mathrm{sm}})\cong L(\xi)\otimes_L(\Ind_B^G\delta_{\mathrm{sm}})^{\mathrm{sm}},
\]
其中局部代数表示由不可约代数表示 $L(\xi)$ 与光滑特征标 $\delta_{\mathrm{sm}}$ 的光滑抛物诱导的张量积给出。
\end{lem}

回忆 BGG 范畴有内蕴的对偶 $M\mapsto M^\vee$（通常称为 BGG 对偶），大致如下给出：把 $M$ 映到其（有限维）$\mathfrak t$-特征空间之对偶的直和，并利用 $\mathfrak g$ 上的主对合（main involution），以便从左模过渡到右模，并得到局部 $U(\mathfrak b_L)$-有限、而非 $U(\overline{\mathfrak b}_L)$-有限的 $U(\mathfrak g_L)$-模；其中 $\overline{\mathfrak b}$ 表示反 Borel 子群 $\overline{B}$ 的李代数。对 $\mathcal O(U(\mathfrak g_L))$ 中的单模 $M$ 有 $M^\vee\cong M$；但 $(-)^\vee$ 反变，故该函子交换子对象与商对象。设 $\delta=z^\xi\delta_{\mathrm{sm}}$ 是局部代数特征标，光滑部分为 $\delta_{\mathrm{sm}}$，代数部分为 $\xi$。于是定义表示
\begin{equation}\tag{6.2.20}\label{eq:a4-6-2-20}
\mathcal F_B^G(\delta)=F_B^G(M(-\xi)^\vee,\delta_{\mathrm{sm}}).
\end{equation}

\begin{rem}{6.2.21}\label{rem:a4-6-2-21}
由 [OS15] 可知，$\Ind_B^G(\delta)^{\mathrm{an}}$ 与 $\mathcal F_B^G(\delta)$ 有相同的不可约子商（重数也相同）。此外注意：引理 6.2.19 连同典范满射 $M(-\xi)^\vee\to L(-\xi)$ 的存在，蕴含对支配的 $\xi$，表示 $\mathcal F_B^G(z^\xi\delta_{\mathrm{sm}})$ 以局部代数表示 $L(\xi)\otimes_L(\Ind_B^G\delta_{\mathrm{sm}})^{\mathrm{sm}}$ 为商；而该局部代数表示在 $\Ind_B^G(z^\xi\delta_{\mathrm{sm}})^{\mathrm{an}}$ 中作为子表示出现。
\end{rem}

我们预期 Orlik--Strauch 函子可用来定义抛物诱导的一个变体（其导出版本当然取决于范畴 $D(\mathrm{an}\,T)$ 与 $D(\mathrm{an}\,G)$ 的定义）。

\medskip\noindent\textbf{期望 6.2.22.}\label{exp:a4-6-2-22}
存在自然函子
\[
\mathcal F_B^G:\Rep_L^{\mathrm{an}}T\lra\Rep_L^{\mathrm{an}}G,\qquad
\mathcal F_B^G:D(\mathrm{an}\,T)\lra D(\mathrm{an}\,G),
\]
使得对局部代数特征标 $\delta$，表示 $\mathcal F_B^G(\delta)$ 由 \eqref{eq:a4-6-2-20} 给出。

\begin{rem}{6.2.23}\label{rem:a4-6-2-23}
可以合理猜测：存在 $D(\mathrm{an}\,G)$ 与 $D(\mathrm{an}\,T)$ 的内蕴对偶 $\mathbf D_G$ 与 $\mathbf D_T$（类似于 6.1 节讨论的对偶），使得（差一个平移）
\[
\mathcal F_B^G(-)=\mathbf D_G\circ\Ind_B^G(-)^{\mathrm{an}}\circ\mathbf D_T.
\]
$\mathbf D_G$ 在光滑表示上的限制应由 Zelevinsky 对合给出；上述期望提示 $\mathbf D_G$ 应与假设性函子 $\mathfrak A_G^{\mathrm{rig}}$ 下的 Serre 对偶相关，比较注 6.2.7。对主级数表示，Strauch 与 Wu 在 [SW25] 中定义并讨论了这样的对偶。
\end{rem}

\ysrc{265}
\begin{rem}{6.2.24}\label{rem:a4-6-2-24}
Juan Esteban Rodríguez Camargo 向我们解释：利用 Clausen--Scholze 发展的解析叠理论，可以在 $(\mathfrak g,B)$-模的范畴上定义函子 $\mathcal F_B^G$ 的一个推广。该推广以及解析叠上的六函子形式化，可用来构造期望 6.2.22 所预期的函子，并证明注 6.2.23 中预期的与对偶的相容性。由于我们不打算从解析叠的观点讨论该理论（原因很多，其中包括我们本人在处理解析叠方面的不足），我们仍采用上述期望中的表述。
\end{rem}

借助猜想的函子 $\mathcal F_B^G$，可以表述期望 6.2.16 的一个变体。该期望的动机来自下文定理 6.2.28 中 Bezrukavnikov 函子的性质。

\medskip\noindent\textbf{期望 6.2.25.}\label{exp:a4-6-2-25}
存在同构
\[
\mathfrak A_G^{\mathrm{rig}}\bigl(\mathcal F_B^G(\pi)\bigr)\cong
\tilde\beta_{\hat B,*}\tilde\alpha_{\hat B}^{!}
\bigl(\mathfrak A_T^{\mathrm{rig}}(s_G(\pi)[-[F:\Q_p]\rho'])\bigr),
\]
关于 $\pi$ 函子性，只要 $\pi$ 是使 $\mathcal F_B^G(\pi)$ 落在 $D^b_{\mathrm{f.p.}}(\mathrm{an}\,G)$ 中、且 $\tilde\alpha_{\hat B}^{!}(\mathfrak A_T^{\mathrm{rig}}(\pi))$ 的支撑在 $\mathcal X_d$ 上 proper 的表示。

\begin{rem}{6.2.26}\label{rem:a4-6-2-26}
考虑 $F=\Q_p$、$d=2$ 的情形，并回忆注 6.2.10：我们预期 $\mathfrak A_G^{\mathrm{rig}}$ 形如
\[
\pi\mapsto\pi\,\widehat\otimes^{L}_{\mathcal D(G)}\,\mathcal L_\infty,
\]
其中 $\mathcal L_\infty$ 是 $\mathcal X_2$ 上的一族 $\mathcal D(G)$-模。对由 $(\varphi,\Gamma)$-模 $D$ 代表的点 $x\in\mathcal X_2(L)$，我们预期
\[
\mathcal L_\infty\otimes k(x)=(\Pi(D)^{\mathrm{an}})'.
\]
这里若 $D$ 平展，则 $\Pi(D)^{\mathrm{an}}$ 表示通过 $\mathrm{GL}_2(\Q_p)$ 的通常 $p$-进局部 Langlands 对应与 $D$ 对应的表示中的局部解析向量（比较 3.2.5 节的讨论）；若 $D$ 非平展，则采用 Colmez 对该对应的推广 [Col16]。给定特征标 $\delta:T\to L$，我们进而预期
\begin{align*}
H^0\bigl(\mathfrak A_G^{\mathrm{rig}}(\Ind_B^G(\delta))\bigr)\otimes k(x)
&\cong\bigl(\Hom_G(\Ind_B^G(\delta),\Pi(D)^{\mathrm{an}})\bigr)',\\
H^0\bigl(\mathfrak A_G^{\mathrm{rig}}(\mathcal F_B^G(\delta))\bigr)\otimes k(x)
&\cong\bigl(\Hom_G(\mathcal F_B^G(\delta),\Pi(D)^{\mathrm{an}})\bigr)'.
\end{align*}
右端向量空间的维数可以显式定出，因为 Colmez 已对局部解析表示 $\Pi(D)^{\mathrm{an}}$ 给出精确描述（见 [Col14, Col16]）；特别地，除非 $D$ 三角化，这些维数为零。利用期望 6.2.16 与 6.2.25，这些纤维维数的计算与定理 5.3.22 (v) 中纤维维数的计算恰好吻合。
\end{rem}

\medskip\noindent\textbf{6.2.27. BGG 范畴上 Bezrukavnikov 的函子.}
我们进一步阐述 Orlik--Strauch 函子及其与抛物诱导的相容性。由于在伽罗瓦侧抛物诱导是用 $\mathcal X_B$ 的紧化 $\overline{\mathcal X}_B$ 来定义的，而 $\overline{\mathcal X}_B$ 在若干正则点处的几何可以显式控制（定理 5.3.40），我们可以猜想式地把函子 $\mathfrak A_G^{\mathrm{rig}}$ 在某些 Orlik--Strauch 表示上的取值与 Bezrukavnikov [Bez16] 在 BGG 范畴上构造的函子联系起来。

\ysrc{266}
固定如上 $\underline\lambda$ 与 $\xi$。由于 $\mathcal O(U(\mathfrak g_L))_{(\xi)}$ 的每个对象都是形如 $M(w\cdot\xi)$（$w\in W$）的 Verma 模的（子）商之扩张，关于函子 $\mathfrak A_G^{\mathrm{rig}}$ 的诸预期蕴含：给定 $\delta_{\mathrm{sm}}$，应有函子
\[
\mathfrak A_G^{\mathrm{rig}}\circ F_B^G\bigl((-)^\vee,\delta_{\mathrm{sm}}\bigr):
\mathcal O(U(\mathfrak g_L))_{(\xi)}\lra D^b_{\mathrm{coh}}(\mathcal X_d),
\]
使该函子像中的层都支撑在定理 5.3.40 之前定义的闭子叠 $\mathcal X_{d,(\underline\lambda,\delta_{\mathrm{sm}})-\mathrm{tri}}$ 上。我们在该函子中额外插入对偶 $(-)^\vee$，以使函子协变。利用定理 5.3.40 中叠 $\widehat{\mathcal X}_{d,(\underline\lambda,\delta_{\mathrm{sm}})-\mathrm{tri}}$ 的局部描述，可以构造该函子，依据是 Bezrukavnikov 的工作。

本节余下部分记
\[
H=(\operatorname{Res}_{F/\Q_p}\mathrm{GL}_d)_L=\prod_{\sigma:F\hookrightarrow L}\mathrm{GL}_d.
\]
注意作为 $L$ 上的李代数有 $\mathfrak g_L\cong\operatorname{Lie}H=\mathfrak h$，特别地 $\mathcal O(U(\mathfrak g_L))$ 与 $\mathcal O(U(\mathfrak h))$ 同构。回忆 5.3.23 节（那里用 $G$ 代替 $H$）：用
\[
\overline{X}_w\subseteq\tilde{\mathfrak h}\times_{\mathfrak h}\tilde{\mathfrak h}
\]
表示支配 Schubert 胞腔 $H(1,w)\subseteq H/B_H\times H/B_H$ 之闭包的不可约分支，其中 $B_H\subseteq H$ 是 $\mathrm{GL}_d$ 中上三角矩阵的 Weil 限制。此外注意有典范投影
\[
\omega:\tilde{\mathfrak h}\times_{\mathfrak h}\tilde{\mathfrak h}
\xrightarrow{\ \mathrm{pr}_1\ }\tilde{\mathfrak h}\lra\widehat{\mathfrak t}
\]
（其中 $\widehat{\mathfrak t}$ 是 $\mathfrak t_L$ 的对偶李代数）；用 $\overline{X}_{w,0}$ 表示 $\overline{X}_w$ 与 $0$ 的原像的交。最后注意零截面给出一个副本
\[
(\operatorname{Res}_{F/\Q_p}\mathrm{GL}_d/P_{\underline\lambda})_L\times
(\operatorname{Res}_{F/\Q_p}\mathrm{GL}_d/P_{\underline\lambda})_L
=H/B_H\times H/B_H\subseteq\overline{X}_{w_0,0}\subseteq\overline{X}_{w_0}
\subseteq\tilde{\mathfrak h}\times_{\mathfrak h}\tilde{\mathfrak h}.
\]
这里 $P_{\underline\lambda}$ 表示 5.2 节中与 $\underline\lambda$ 相伴的抛物子群；由于假定 $\underline\lambda$ 正则，它与 Borel 子群 $B$ 一致。下述函子由 Bezrukavnikov 在 [Bez16, 11.4] 中构造，另见 [BK, Theorem 13.0.1]。

\begin{thm}{6.2.28 (Bezrukavnikov)}\label{thm:a4-6-2-28}
存在正合函子
\[
F_\xi:\mathcal O(U(\mathfrak g_L))_{(\xi)}\lra
\Coh^H\bigl((\tilde{\mathfrak h}\times_{\mathfrak h}\tilde{\mathfrak h})^{\wedge}_0\bigr),
\]
使得
\begin{enumerate}[label=(\roman*),leftmargin=2.2em,itemsep=0.15em]
\item 对 $w\in W$，Verma 模 $M(ww_0\cdot\xi)$ 的像是 $\overline{X}_w$ 的对偶层：
\[
F_\xi(M(ww_0\cdot\xi))=\omega_{\overline{X}_w}.
\]
\item 对 $w\in W$，对偶 Verma 模 $M(ww_0\cdot\xi)^\vee$ 的像是 $\overline{X}_w$ 的结构层：
\[
F_\xi(M(ww_0\cdot\xi)^\vee)=\mathcal O_{\overline{X}_w}.
\]
\item 不可约代数表示 $L(\xi)$\ysrc{267} 的像是
\[
(\operatorname{Res}_{F/\Q_p}\mathrm{GL}_d/P_{\underline\lambda})_L\times
(\operatorname{Res}_{F/\Q_p}\mathrm{GL}_d/P_{\underline\lambda})_L
=H/B_H\times H/B_H\subseteq\overline{X}_{w_0,0}
\]
上与特征标 $((-\rho')_\sigma,(-\rho')_\sigma)$ 相伴的线丛（$\rho'$ 的定义比较注 6.2.8）。
\end{enumerate}
\end{thm}

我们预期：利用定理 5.3.40 中通过映射
\[
f_{\underline\lambda,\delta_{\mathrm{sm}}}:
\widehat{\mathcal X}_{d,(\underline\lambda,\delta_{\mathrm{sm}})-\mathrm{tri}}
\lra(\tilde{\mathfrak h}\times_{\mathfrak h}\tilde{\mathfrak h})^{\wedge}_0
\]
给出的叠 $\widehat{\mathcal X}_{d,(\underline\lambda,\delta_{\mathrm{sm}})-\mathrm{tri}}$ 的局部描述，Bezrukavnikov 函子的拉回经 Orlik--Strauch 函子给出猜想函子 $\mathfrak A_G^{\mathrm{rig}}$ 的限制：

\medskip\noindent\textbf{期望 6.2.29.}\label{exp:a4-6-2-29}
设 $\varphi_1,\dots,\varphi_d\in L^\times$ 满足 $i\neq j$ 时 $\varphi_i/\varphi_j\neq1,q$，并设 $\delta_{\mathrm{sm}}$ 如 (5.3.17)。则函子
\[
\mathfrak A_G^{\mathrm{rig}}\circ\mathcal F_B^G(-,\delta_{\mathrm{sm}}):
\mathcal O(U(\mathfrak g_L))_{(\xi)}\lra D^b_{\mathrm{coh}}(\mathcal X_d)
\]
经过 $\widehat{\mathcal X}_{d,(\underline\lambda,\delta_{\mathrm{sm}})-\mathrm{tri}}$ 分解，且由复合
\[
\bigl(f^{*}_{\underline\lambda,\delta_{\mathrm{sm}}}\circ F_\xi(-)\bigr)\bigl([F:\Q_p]\rho'\bigr)
\]
给出。

\begin{rem}{6.2.30}\label{rem:a4-6-2-30}
此处用 $\rho'$ 作扭在 $\widehat{\mathcal X}_{d,(\underline\lambda,\delta_{\mathrm{sm}})-\mathrm{tri}}$ 上有意义，比较定理 5.3.40 之后的讨论。注意：利用引理 6.2.19，以及 $f_{\underline\lambda,\delta_{\mathrm{sm}}}$ 形式光滑、其相对对偶层为与 $-2\rho$ 相伴的线丛这一事实，该期望与 (6.2.5) 中用 $(-\rho')_\sigma$ 作的扭、以及期望 6.2.16 与 6.2.25 中的扭相容。
\end{rem}

\medskip\noindent\textbf{6.2.31. 三角化的 Breuil--Mézard 猜想.}
我们描述与抛物诱导的相容性、以及与猜想 6.2.4 中光滑范畴 Langlands 的相容性的一个影子。这个影子可视为 Breuil--Mézard 猜想（猜想 6.1.39）在局部解析主级数表示上的版本，其中涉及的是 Grothendieck 群而非闭链群（比较 6.1.38 节的讨论）。在某些正则情形，涉及闭链的这一猜想的版本已在 [BHS17b] 中证明，见下文定理 6.2.37。

如上，记 $G=\mathrm{GL}_d(F)$，$T\subseteq B\subset G$ 分别为对角环面与上三角 Borel 子群。仍假定 $L$ 足够大。注意可把 $G$ 看作 $H=\operatorname{Res}_{F/\Q_p}\mathrm{GL}_d$ 的 $\Q_p$-值点群；类似地，$T$ 与 $B$ 是对角环面与上三角 Borel 子群的 Weil 限制 $T_H$ 与 $T_B$ 的 $\Q_p$-值点。固定正则 Hodge--Tate 权 $\underline\lambda$，并照例用 $\xi$ 表示相应的平移特征标（它是 $\bar F$ 上定义的 $T_H$ 的特征标）。把 $\xi$ 看作（代数）特征标 $z^\xi:T\to L^\times$（比较 (5.3.18)），并固定光滑特征标 $\delta_{\mathrm{sm}}:T\to L^\times$。

记 $W\cong\prod_{K\hookrightarrow L}\mathcal S_n$ 为 $H$ 的绝对 Weyl 群，$W_0\cong\mathcal S_n\subseteq W$ 为 $H$ 的相对 Weyl 群。\tn{底本此处作 $\prod_{K\hookrightarrow L}\mathcal S_n$ 与 $\mathcal S_n$；按上下文（$H=(\operatorname{Res}_{F/\Q_p}\mathrm{GL}_d)_L$、$G=\mathrm{GL}_d(F)$）疑应为 $\prod_{\sigma:F\hookrightarrow L}\mathcal S_d$ 与 $\mathcal S_d$，译文照录。} 则对每个 $w\in W$ 与 $w'\in W_0$ 都得到良定义的特征标 $z^{w\cdot\xi}$ 与 $w'\delta_{\mathrm{sm}}$。

以 $K_0(\underline\lambda,\delta_{\mathrm{sm}})$ 表示由局部解析诱导表示 $\Ind_B^G(z^{w\cdot\xi}w'\delta_{\mathrm{sm}})$（$w\in W$，$w'\in W_0$）的不可约成分生成的自由阿贝尔群。

\ysrc{268}
\begin{rem}{6.2.32}\label{rem:a4-6-2-32}
记号 $K_0$ 应使读者想起：这个自由阿贝尔群是有限长局部解析表示范畴的 Grothendieck 群，而这些表示的不可约子商都出现在上述局部解析主级数表示中。
\end{rem}

在伽罗瓦侧，回忆定理 5.3.40 之前定义的叠 $\mathcal X_{d,(\underline\lambda,\delta_{\mathrm{sm}})-\mathrm{tri}}$。考虑并
\[
\mathcal X_{d,[\underline\lambda,\delta_{\mathrm{sm}}]-\mathrm{tri}}
=\bigcup_{w'\in W_0}\mathcal X_{d,(\underline\lambda,w'\delta_{\mathrm{sm}})-\mathrm{tri}},
\]
以及该叠上凝聚层的 Grothendieck 群 $K_0(\Coh(\mathcal X_{d,[\underline\lambda,\delta_{\mathrm{sm}}]-\mathrm{tri}}))$。注意 $\mathcal X_{d,[\underline\lambda,\delta_{\mathrm{sm}}]-\mathrm{tri}}$ 恰是如下 $(\varphi,\Gamma)$-模之集：它容许一个滤过，其分次片为 $\mathcal R_F(z^{w\underline\lambda}w'\delta_{\mathrm{sm}})$，对某个 $w\in W$ 与 $w'\in W_0$。

光滑特征标 $\delta_{\mathrm{sm}}$ 定义了 $G$ 的光滑表示范畴的一个 Bernstein 分量 $\Omega=[T,\delta_{\mathrm{sm}}]$；用 $\tau$ 表示使 $\Omega=\Omega_\tau$ 的惯性型。如 5.2.11 节所回顾，分量 $\Omega_\tau$ 定义了秩 $d$ 的 Weil--Deligne 表示之叠、以及秩 $d$、Hodge--Tate 权为 $\underline\lambda$ 的滤过 Weil--Deligne 表示之叠的连通分支
\[
\mathrm{WD}_{d,F,\tau}\subseteq\mathrm{WD}_{d,F},\qquad
\mathrm{Fil}_{\underline\lambda}\mathrm{WD}_{d,F,\tau}\subseteq
\mathrm{Fil}_{\underline\lambda}\mathrm{WD}_{d,F}.
\]
连通分支 $\mathrm{WD}_{d,F,\tau}$ 带有典范映射
\begin{equation}\tag{6.2.33}\label{eq:a4-6-2-33}
\mathrm{WD}_{d,F,\tau}\lra\Spec\mathcal Z_{\Omega_\tau},
\end{equation}
见 (5.2.15)；而 $\delta_{\mathrm{sm}}$ 定义了点 $[\delta_{\mathrm{sm}}]\in\Spec\mathcal Z_{\Omega_\tau}$。用
\[
\mathrm{WD}_{d,F}(\delta_{\mathrm{sm}})\subseteq\mathrm{WD}_{d,F,\tau}
\]
表示 $[\delta_{\mathrm{sm}}]$ 在 \eqref{eq:a4-6-2-33} 下的原像。

\begin{rem}{6.2.34}\label{rem:a4-6-2-34}
用（光滑）范畴局部 Langlands 来说，这个闭子叠可刻画为下述凝聚层支撑之并：这些层是 $\mathfrak A_G^{\mathrm{sm}}(\pi)$，附着于出现在 $\iota_B^G(\delta_{\mathrm{sm}})$ 中的不可约光滑表示 $\pi$，其中 $\iota_B^G$ 表示归一化抛物诱导。注意无需考虑所有可能的 $\iota_B^G(w'\delta_{\mathrm{sm}})$（$w'\in W_0$），因为 $\iota_B^G(\delta)_{\mathrm{sm}}$ 与 $\iota_B^G(w'\delta)_{\mathrm{sm}}$ 有相同的 Jordan--Hölder 因子。
\end{rem}

\begin{lem}{6.2.35}\label{lem:a4-6-2-35}
引理 5.2.3 的同构 $\mathrm{Fil}_{\underline\lambda}\mathrm{WD}_{d,F,\tau}\cong\mathrm{Fil}_{\underline\lambda}\mathrm{Mod}_{d,\varphi,N,\tau}$ 诱导闭嵌入
\[
\mathrm{Fil}_{\underline\lambda}\mathrm{WD}_{d,F,\tau}(\delta_{\mathrm{sm}})^{\mathrm{an}}
\subseteq\mathcal X_{d,(\underline\lambda,\delta_{\mathrm{sm}})-\mathrm{tri}}.
\]
\end{lem}

\begin{proof}
由右端是以 $z^w\underline\lambda w'\delta_{\mathrm{sm}}$ 为参数的三角化 $(\varphi,\Gamma)$-模之叠这一事实立得。
\end{proof}

下述猜想中的映射应理解为 $\mathfrak A_G^{\mathrm{rig}}$ 在 Grothendieck 群上诱导的映射。

\ysrc{269}
\begin{conj}{6.2.36}\label{conj:a4-6-2-36}
存在唯一的单射群同态
\[
\mathbf a_{\underline\lambda,\delta_{\mathrm{sm}}}:
K_0(\underline\lambda,\delta_{\mathrm{sm}})\lra
K_0\bigl(\Coh(\mathcal X_{d,[\underline\lambda,\delta_{\mathrm{sm}}]-\mathrm{tri}})\bigr),
\]
使得
\begin{enumerate}[label=(\roman*),leftmargin=2.2em,itemsep=0.2em]
\item 对 $w\in W$ 与 $w'\in W_0$ 有
\begin{align*}
&\mathbf a_{\underline\lambda,\delta_{\mathrm{sm}}}\bigl([\Ind_B^G(z^{ww_0\cdot\xi}w'\delta_{\mathrm{sm}})^{\mathrm{an}}]\bigr)\\
&\qquad=\Bigl[\tilde\beta_{\hat B,*}\circ\tilde\alpha_{\hat B}^{*}
\Bigl(\mathfrak A_M^{\mathrm{rig}}\Bigl(s_G\bigl(z^{ww_0\cdot\xi}w'\delta_{\mathrm{sm}}\bigr)\\
&\qquad\qquad\qquad\qquad\qquad\qquad
\bigl([F:\Q_p]\rho'\bigr)\Bigr)\Bigr)\Bigr].
\end{align*}
\item 对区段 $\Rep^{\mathrm{sm}}_{\Omega_\tau}G$ 中作为 $\iota_B^G\delta_{\mathrm{sm}}$ 的成分出现的不可约光滑表示 $\pi$，有
\[
\mathbf a_{\underline\lambda,\delta_{\mathrm{sm}}}([\pi\otimes L(\xi)])
=\bigl[\bigl(\mathrm{pr}_{\mathrm{WD}}^{*}(\mathfrak A_G^{\mathrm{sm}}(\pi))\bigr)(-\rho')_\sigma\bigr].
\]
\end{enumerate}
\end{conj}

当然，公式 (i) 与 (ii) 应与猜想 6.2.4 中的性质 (1) 与 (2) 对照。该猜想的一个较弱版本在更严格的假设下已于 [BHS17b, Theorem 4.3.8] 中证明。事实上，loc. cit. 的证明从形变环的谱整体化到叠 $\mathcal X_{d,[\underline\lambda,\delta_{\mathrm{sm}}]-\mathrm{tri}}$ 上。更确切地说，以 $\mathrm{CH}^0(\mathcal X_{d,[\underline\lambda,\delta_{\mathrm{sm}}]-\mathrm{tri}})$ 表示 $\mathcal X_{d,[\underline\lambda,\delta_{\mathrm{sm}}]-\mathrm{tri}}$ 的不可约分支上的自由阿贝尔群。把凝聚层映到其支撑（并忘掉非极大维数的分支）定义了态射
\[
\mathrm{supp}:K_0\bigl(\Coh(\mathcal X_{d,[\underline\lambda,\delta_{\mathrm{sm}}]-\mathrm{tri}})\bigr)
\lra\mathrm{CH}^0(\mathcal X_{d,[\underline\lambda,\delta_{\mathrm{sm}}]-\mathrm{tri}}).
\]
于是 [BHS17b, Theorem 4.3.8] 的证明方法可用来证明下述定理。

\begin{thm}{6.2.37}\label{thm:a4-6-2-37}
设 $\delta_{\mathrm{sm}}=(\mathrm{unr}\varphi_1,\dots,\mathrm{unr}\varphi_n)$ 是非分歧特征标，且 $i\neq j$ 时 $\varphi_i/\varphi_j\notin\{1,q\}$。则存在唯一的同态
\[
\mathbf a'_{\underline\lambda,\delta_{\mathrm{sm}}}:
K_0(\underline\lambda,\delta_{\mathrm{sm}})\lra
\mathrm{CH}^0(\mathcal X_{d,[\underline\lambda,\delta_{\mathrm{sm}}]-\mathrm{tri}}),
\]
使得把 $\mathbf a_{\underline\lambda,\delta_{\mathrm{sm}}}$ 换成 $\mathbf a'_{\underline\lambda,\delta_{\mathrm{sm}}}$、并把等式右端与支撑映射 $\mathrm{supp}$ 复合之后，(i) 与 (ii) 的条件成立。
\end{thm}

\begin{rem}{6.2.38}\label{rem:a4-6-2-38}
（a）事实上，在猜想 6.2.36 中还可要求更多相容性，即使用比 $T$ 更一般的 Levi 子群的抛物诱导。但这会使猜想的表述稍显繁复。

（b）猜想 6.2.4 相当难以攻破，因为尚不存在满足诸约束的函子候选；相比之下，猜想 6.2.36 明确得多：如同 Breuil--Mézard 猜想，群同态 $\mathbf a_{\underline\lambda,\delta_{\mathrm{sm}}}$ 由猜想中的条件 (i) 与 (ii) 唯一确定，猜想更多是检验关系，而非构造映射（不过此处与 Breuil--Mézard 猜想相反，只有有限多条关系需要检验）。例如在定理 6.2.37 的附加假设下，Bezrukavnikov 的函子 6.2.28 与期望 6.2.29 已经确定了 $\mathbf a_{\underline\lambda,\delta_{\mathrm{sm}}}$（只要能够验证相应的关系）。

（c）猜想 6.2.36 涉及 Grothendieck 群的版本比其涉及闭链的版本精细得多。例如在 $d=2$ 的情形，若 $\pi$ 是平凡表示 $\mathbf 1$ 与（光滑）Steinberg 表示 $\mathrm{St}$ 的非分裂扩张这一（抛物诱导的）光滑表示，则 $\mathbf a_{\underline\lambda,\delta_{\mathrm{sm}}}(L(\xi)\otimes\pi)$ 底层的闭链平凡，而 $\mathbf a_{\underline\lambda,\delta_{\mathrm{sm}}}(L(\xi)\otimes\pi)$ 在凝聚层 Grothendieck 群中的类非平凡（比较如 [Hel23, Remark 4.43]）。更确切地说，函子 $\mathfrak A_G^{\mathrm{sm}}$ 把 $\mathrm{St}$ 映到不可约分支 $X_1\subseteq\mathrm{WD}_{2,F,\mathbf 1}$（秩 2 的 $(\varphi,N)$-模之叠，其中单值 $N$ 一般非平凡）的结构层 $\mathcal O_{X_1}$（在次数 0），\ysrc{270}而把 $\mathbf 1$ 映到 $X_1$ 上平移至上同调次数 $-1$ 的线丛 $\mathcal L$。线丛 $\mathcal L$ 与结构层 $\mathcal O_{X_1}$ 作为叠 $X_1$ 上的凝聚层不同构，但当然有相同的底层闭链。在 $\mathfrak A_G^{\mathrm{sm}}(\pi)$ 所定义闭链的公式中，这两个闭链相消，因为它们所在的上同调次数奇偶性不同。另一方面，$\mathfrak A_G^{\mathrm{sm}}(\pi)$ 在凝聚层 Grothendieck 群中的类仍非平凡。

（d）$\mathbf a'_{\underline\lambda,\delta_{\mathrm{sm}}}$ 的构造位于 9.6.33 节特征簇上伴点（companion point）构造的核心。粗略地说，可与 Breuil--Mézard 猜想作如下类比：证明 Breuil--Mézard 猜想蕴含（在某些良好情形）可以构造具有指定伽罗瓦表示的自守形式；类似地，定理 6.2.37 的局部解析 Breuil--Mézard 重数公式蕴含可以构造（超收敛有限斜率）具有指定伽罗瓦表示的 $p$-进自守形式。
\end{rem}

\section{范畴 $p$-进 Langlands 的若干已知情形}

\subsection{$\mathrm{GL}_1$}

不出所料，猜想 6.1.15 的 $\mathrm{GL}_1$ 情形是局部类域论的推论。此外，在此情形我们可以描述函子的本质像（essential image），而且无需过渡到导出层面，因为该函子来自阿贝尔范畴的等价。关键是：由局部类域论有同构 $W_F^{\mathrm{ab}}\isom F^\times=\mathrm{GL}_1(F)$；而由 [EG23, Rem. 7.2.18]，叠 $\mathcal X_1$ 可描述为 $W_F$ 的 1 维连续表示、从而 $W_F^{\mathrm{ab}}\isom\mathrm{GL}_1(F)$ 的表示的模叠。

更确切地说，选取单值化元 $\varpi_F\in F$ 后可写 $F^\times=\mathcal O_F^\times\times\varpi_F^{\Z}$，从而
\[
\mathcal O[[F^\times]]=\mathcal O[[\mathcal O_F^\times]]\widehat\otimes_{\mathcal O}\mathcal O[X,X^{-1}]^{\wedge},
\]
其中 $\mathcal O[X,X^{-1}]^{\wedge}$ 是 $\mathcal O[X,X^{-1}]$ 的 $p$-进完备化，张量积也取 $p$-进完备化。于是由 [EG23, Prop. 7.2.17]（另见 [Pha24, Cor. 1.2] 中 Dat Pham 的另一个证明，其思路与下文提到的 Kedlaya--Pottharst--Xiao 方法类似），有同构
\begin{equation}\tag{7.1.1}\label{eq:a4-7-1-1}
\bigl[\bigl(\Spf\mathcal O[[F^\times]]\bigr)/\widehat{\mathbb G}_m\bigr]\isom\mathcal X_1
\end{equation}
（其中 $\widehat{\mathbb G}_m=\Spf\mathcal O[X,X^{-1}]^{\wedge}$ 是 $\mathbb G_m$ 的 $p$-进完备化；构造商叠时，$\widehat{\mathbb G}_m$-作用取为平凡）。

由 \eqref{eq:a4-7-1-1} 立得（阿贝尔）范畴 $\mathcal O[[F^\times]]\text{-}\Mod$ 与 $\mathcal X_1$ 上 $\widehat{\mathbb G}_m$-作用平凡的拟凝聚层范畴之间的等价。表述该等价的另一方式是：结构层 $\mathcal O_{\mathcal X_1}$ 自然地是 $\mathcal O[[F^\times]]$-模，而我们的函子由
\begin{equation}\tag{7.1.2}\label{eq:a4-7-1-2}
\pi\mapsto\mathcal O_{\mathcal X_1}\otimes_{\mathcal O[[F^\times]]}\pi
\end{equation}
给出。换言之，在 $d=1$ 的情形，注 6.1.24 中的层 $\mathcal L_\infty$ 就是 $\mathcal O_{\mathcal X_1}$。

叠 $\mathcal X_1$ 还有平行的描述，以及用 $F^\times$ 的局部解析表示对其上凝聚层所作的解释。由 Kedlaya--Pottharst--Xiao 的结果 [KPX14, Thm. 6.2.14]，对仿射代数 $A$，$R_{F,A}$ 上每个秩一 $(\varphi,\Gamma)$-模都形如 $R_{F,A}(\delta)\otimes\mathcal L$，其中 $\delta:F^\times\to A^\times$ 是连续特征标，$\mathcal L$ 是 $\Sp A$ 上的线丛。

\ysrc{271}
由此
\begin{equation}\tag{7.1.3}\label{eq:a4-7-1-3}
\mathcal X_1=\mathcal T/\mathbb G_m
=\bigl((\Spf\mathcal O[[\mathcal O_F^\times]])^{\mathrm{rig}}\times\mathbb G_m\bigr)/\mathbb G_m,
\end{equation}
其中 $\mathcal T=\Hom_{\mathrm{cont}}(F^\times,\mathbb G_m(-))$ 是 $F^\times$ 的连续特征标构成的刚性解析空间，配有刚性解析乘法群 $\mathbb G_m$ 的平凡作用。\tn{注意我们通常用 $\mathbb G_m$ 表示与概形 $\mathbb G_m$ 相伴的刚性解析空间。需要把 $\mathbb G_m$ 看作概形时，我们总会明确说明，故不应引起混淆。} 如上，单值化元的选取使我们可写 $\mathcal T=\mathcal W\times\mathbb G_m$，其中
\[
\mathcal W=\Hom_{\mathrm{cont}}(\mathcal O_F^\times,\mathbb G_m(-))
=(\Spf\mathcal O[[\mathcal O_F^\times]])^{\mathrm{rig}}
\]
是 $\mathcal O_F^\times$ 的连续特征标空间。特别地，我们看到有典范映射
\[
\mathcal X_{1,\eta}^{\mathrm{rig}}
=(\mathcal W\times\widehat{\mathbb G}_m^{\mathrm{rig}})/\widehat{\mathbb G}_m^{\mathrm{rig}}
\lra(\mathcal W\times\mathbb G_m)/\mathbb G_m=\mathcal X_1.
\]
该映射的像是 $\mathcal X_1$ 的开子叠；注意该映射不可表。其非平凡纤维是零维的，由商叠 $\mathbb G_m/\widehat{\mathbb G}_m^{\mathrm{rig}}$ 给出。

有同构
\[
D(\mathcal O_F^\times)\cong\Gamma(\mathcal W,\mathcal O_{\mathcal W}),
\]
即 $\mathcal O_F^\times$ 的局部解析分布代数与 $\mathcal O_F^\times$ 的特征标空间 $\mathcal W$ 上结构层的整体截面同构。注意 $D(F^\times)=D(\mathcal O_F^\times)[T^{\pm1}]$ 是 $D(\mathcal O_F^\times)$ 上的 Laurent 多项式环，即 $D(\mathcal O_F^\times)$ 上代数 $\mathbb G_m$ 的函数环；而
\[
\Gamma(\mathcal X_1,\mathcal O_{\mathcal X_1})=\Gamma(\mathcal T,\mathcal O_{\mathcal T})
\cong\Gamma(\mathcal W,\mathcal O_{\mathcal W})\widehat\otimes_L\Gamma(\mathbb G_m,\mathcal O_{\mathbb G_m})
\]
是 $D(\mathcal O_F^\times)$ 上刚性解析 $\mathbb G_m$ 的函数空间。

于是，如 \eqref{eq:a4-7-1-2}，我们想定义函子
\begin{equation}\tag{7.1.4}\label{eq:a4-7-1-4}
\pi\mapsto\mathcal O_{\mathcal X_1}\widehat\otimes^{L}_{D(F^\times)}\pi,
\end{equation}
其中，限制到 $\mathcal X_1$ 中的仿射子空间之后，符号 $\mathcal O_{\mathcal X_1}\widehat\otimes^{L}_{D(F^\times)}\pi$ 应是导出完备张量积（并且再次强调：至少导出的部分只在固态模的世界中才有意义）。与 $p$ 幂挠表示的情形不同，函子 \eqref{eq:a4-7-1-4} 并非定义在阿贝尔范畴层面；因此在局部解析表示的设定中，范畴 $p$-进 Langlands 对应即使在 $\mathrm{GL}_1$ 情形也只是导出层面的陈述！为把这一点说得更精确，最好考虑 \eqref{eq:a4-7-1-4} 的伴随函子。

为简化讨论，考虑 $F=\Q_p$ 的情形，并进一步限制到 $\Z_p^\times$ 而非 $\Q_p^\times$。这一情形的讨论同样容易化归到 $(\Z_p,+)\cong(1+p\Z_p,\cdot)$ 的情形。

回忆分布代数 $D(\Z_p)$ 同构于 $\Q_p$ 上开单位圆盘 $\mathbf U$ 上解析函数环 $\Gamma(\mathbf U,\mathcal O_{\mathbf U})$。这里用同构 $\Z_p[[\Z_p]]\cong\Z_p[[T]]$，把 $(\Spf\Z_p[[\Z_p]])^{\mathrm{rig}}$（即 $\Z_p$ 的连续特征标空间）与开单位圆盘 $\mathbf U=(\Spf\Z_p[[T]])^{\mathrm{rig}}$ 等同。

\ysrc{272}
刚性解析空间 $\mathbf U$ 是光滑 Stein 空间。对这类空间上的凝聚层，Chiarelotto [Chi90] 构造了紧支撑上同调 $R\Gamma_c(\mathbf U,-)$ 以及 Serre 对偶
\begin{equation}\tag{7.1.5}\label{eq:a4-7-1-5}
H^0(\mathbf U,\mathcal F)'\cong H^1_c(\mathbf U,\mathbf D(\mathcal F)),
\end{equation}
其中 $(-)'$ 表示强对偶（如上），$\mathbf D(\mathcal F)=\mathbf D_{\mathcal O_{\mathbf U}}(\mathcal F)$ 表示凝聚层 $\mathcal F$ 的 Serre 对偶。由函子性，$H^1_c(\mathbf U,\mathbf D(\mathcal F))$（更一般地 $R\Gamma_c(\mathbf U,\mathcal F^\bullet)$）是 $D(\Z_p)$-模（或这种模的复形）。事实上，$H^1_c$ 上诱导的 $\Z_p$-作用是局部解析（且容许）的，因为按 Serre 对偶，它是 $D(\Z_p)$ 上一个余容许模（coadmissible module）的强对偶。函子 \eqref{eq:a4-7-1-2} 现在应定义为函子
\[
\mathcal F\mapsto R\Gamma_c(\mathbf U,\mathcal F)
\]
的右伴随函子，该函子从凝聚层的导出范畴 $D^b_{\mathrm{coh}}(\mathbf U)$ 映到固态 $\Q_p$-向量空间中 $D(\Z_p)$-模的某个导出范畴（Clausen--Scholze 意义下）。

注意函子 $R\Gamma_c(\mathbf U,-)$ 不保持凝聚层的上同调次数：若 $\mathcal F$ 是集中在次数 0、支撑有限的凝聚层，则 $R\Gamma_c(\mathbf U,\mathcal F)$ 与整体截面 $R\Gamma(\mathbf U,\mathcal F)$ 一致，集中在次数 0，且是有限维 $\Q_p$-向量空间。另一方面，$R\Gamma_c(\mathbf U,\mathcal O_{\mathbf U})$ 集中在次数 1，此处它由 $\Z_p$ 上局部解析函数空间 $C^{\mathrm{an}}(\Z_p,\Q_p)$ 给出（因为 $C^{\mathrm{an}}(\Z_p,\Q_p)$ 等同于 $D(\Z_p)=R\Gamma(\mathbf U,\mathcal O_{\mathbf U})$ 的强对偶）。

反之，伴随函子 $-\widehat\otimes^{L}_{D(\Z_p)}\mathcal O_{\mathbf U}$ 把有限维局部解析表示映为次数 0 的凝聚层，而把 $C^{\mathrm{an}}(\Z_p,\Q_p)$ 映为 $\mathcal O_{\mathbf U}[1]$，即平移至次数 $-1$ 的结构层。

\begin{rem}{7.1.6}\label{rem:a4-7-1-6}
若只考虑集中在次数 0 的模或表示，则非导出版本 $\pi\otimes_{D(\Z_p)}\mathcal O_{\mathbf U}$ 也可以（至少对某些 $\pi$）不借助凝聚态与固态世界来计算，它给出 Stein 空间 $\mathbf U$ 上附着于结构层整体截面环 $D(\Z_p)=\Gamma(\mathbf U,\mathcal O_{\mathbf U})$ 上某个模的层。但读者应注意：该完备张量积不与取整体截面交换，即
\[
\Gamma(\mathbf U,\pi\otimes_{D(\Z_p)}\mathcal O_{\mathbf U})
\neq\pi\otimes_{D(\Z_p)}\Gamma(\mathbf U,\mathcal O_{\mathbf U})=\pi.
\]
例如当 $\pi=C^{\mathrm{an}}(\Z_p,\Q_p)$ 时，左端为 0，因为层 $\pi\otimes_{D(\Z_p)}\mathcal O_{\mathbf U}$ 为零：$C^{\mathrm{an}}(\Z_p,\Q_p)$ 与 $\mathbf U$ 的每个开仿射子空间上函数环的完备张量积都消失。然而，如上所述，导出版本 $\pi\widehat\otimes^{L}_{D(\Z_p)}\mathcal O_{\mathbf U}$ 非零；用 $R\Gamma_c(\mathbf U,-)$ 代替 $R\Gamma(\mathbf U,-)$ 便还原出表示 $\pi$。
\end{rem}

\begin{rem}{7.1.7}\label{rem:a4-7-1-7}
利用 Clausen--Scholze 引入的固态凝聚模与解析环 [Sch19]（以及在刚性解析空间情形下由 Andreychev [And21] 进一步研究），确实可以把上述内容精确化。为避免过多涉及凝聚态结构，我们不在此展开。但要注意，应把上述情形与仿射直线的紧化
\[
\iota:\Spa(\Z[x],\Z[x])\lra\Spa(\Z[x],\Z)
\]
相比较。\ysrc{273}
在离散 adic 空间的语境中，见 [Sch19, Theorem 8.1]。在此语境中，导出完备化
\[
\iota^*=-\otimes^{L}_{(\Z[x],\Z)_{\blacksquare}}(\Z[x],\Z[x])_{\blacksquare}:
\mathbf D_{\blacksquare}(\Z[x],\Z)\lra\mathbf D_{\blacksquare}(\Z[x],\Z[x])
\]
有全忠实左伴随 $\iota_!$（紧支撑上同调），且 $\iota^*\circ\iota_!$ 是恒等，见 [Sch19, Observation 8.11]。此外 $\iota_!(\Z[x])=(\Z((x^{-1}))/\Z[x])[-1]$ 位于上同调次数 1（这蕴含导出完备化 $\iota^*$ 把 $\Z((x^{-1}))/\Z[x]$ 映到平移至次数 $-1$ 的结构层），见 [Sch19, Observation 8.12]。
\end{rem}

上述讨论可概括为下述定理，它是 [RJRC25, Theorem B] 的特例。这里用 $D_{qc}(\mathbf U)$ 表示 Clausen--Scholze 意义下 $\mathbf U$ 上固态拟凝聚层的导出范畴。

\begin{thm}{7.1.8}\label{thm:a4-7-1-8}
函子 $R\Gamma_c(\mathbf U,-)$ 诱导范畴等价
\[
\begin{aligned}
D_{qc}(\mathbf U)&\isom D(\mathrm{an}\,\Z_p),\\
D^b_{\mathrm{coh}}(\mathbf U)&\isom D_{\mathrm{f.p.}}(\mathrm{an}\,\Z_p),
\end{aligned}
\]
其逆函子由 $-\widehat\otimes^{L}_{D(\Z_p)}\mathcal O_{\mathbf U}$ 给出。
\end{thm}

在 Banach 情形，函子 \eqref{eq:a4-7-1-2} 是阿贝尔范畴之间函子（的导出函子），而在局部解析设定中函子 \eqref{eq:a4-7-1-4} 只在导出层面才有意义，这似乎令人惊讶。特别地，其中一个函子不能平移上同调次数，另一个却可以。尽管如此，这两个函子在猜想 6.2.4 的意义下相容，以下来说明。

考虑 ind-凝聚层的 pro-系统
\[
\mathcal F^{(n)}_m=F^{(n)}_m\otimes_{\Z_p[[T]]}\mathcal O_{\Spf(\Z_p[[T]])},
\]
定义在 $\Spf(\Z/p^n[[T]])=\varinjlim_m\Spec((\Z/p^n)[T]/T^{p^m})$ 上，其中 $F^{(n)}_m=(\Z/p^n)[\Z_p/p^m\Z_p]$ 表示 $\Z_p$ 在 $\Z_p$ 上取值于 $\Z/p^n$、且在 $p^m\Z_p$-陪集上为常值的函数空间上的表示（这里当然用同构 $\Z_p[[\Z_p]]\cong\Z_p[[T]]$）。在表示层面，这个 pro-ind 系统（取 $\varprojlim_n\varinjlim_m$ 之后）给出 $\Z_p$ 在连续映射空间 $C_{\mathrm{cont}}(\Z_p,\Z_p)$ 上的表示；于是把 $p$ 逆过来并取局部解析向量之后，便得到 $C^{\mathrm{an}}(\Z_p,\Q_p)$。因此 Banach 与解析情形的相容性由下述引理给出。

\begin{lem}{7.1.9}\label{lem:a4-7-1-9}
pro-ind 系统 $(\mathcal F^{(n)}_m)$（集中在次数 0）的一般纤维（generic fiber）$(\mathcal F^{(n)}_m)_\eta$ 是开单位圆盘上平移至次数 $-1$ 的结构层 $\mathcal O_{\mathbf U}[1]$。
\end{lem}

在概述证明之前，须先定义 pro-ind 系统 $(\mathcal F^{(n)}_m)$ 的一般纤维 $(\mathcal F^{(n)}_m)_\eta$。回忆开单位圆盘有容许覆盖
\[
\mathbf U=\bigcup_{k\ge1}\mathbf U_k,
\]
其中 $\mathbf U_k$ 是半径 $p^{-1/k}$ 的闭圆盘，每个闭圆盘都有良好的仿射形式模型 $\mathfrak U_k$。以 $j_k:\mathbf U_k\into\mathbf U$ 表示开嵌入，并记

\ysrc{274}
$\pi_k:\mathfrak U_k\to\Spf\Z_p[[T]]$ 为典范映射。则 $j_k^*((\mathcal F^{(n)}_m)_\eta)$ 按定义是 $\mathfrak U_k$ 上由
\[
\varprojlim_n\varinjlim_m\pi_k^*\mathcal F^{(n)}_m
\]
给出的凝聚层的一般纤维；我们通过在覆盖 $\mathbf U=\bigcup_{k\ge1}\mathbf U_k$ 上粘合这些层来定义 $(\mathcal F^{(n)}_m)_\eta$。

\begin{proof}[引理 7.1.9 的证明概述]
我们概述 $k=1$ 的情形，此时 $\mathfrak U_1=\Spf\Z_p\langle X\rangle$，$X=p^{-1}T$。为简单起见，只证明
\[
\varinjlim_m\bar\pi_1^*\mathcal F^{(1)}_m
\]
是 $\bar{\mathcal U}_1=\Spec\F_p[X]$ 上平移至次数 $-1$ 的结构层，一般情形留作练习。这里
\[
\bar\pi_1:\Spec\F_p[X]\lra\Spf\F_p[[T]]
\]
是 $\pi_1$ 的模 $p$ 纤维，即由 $T\mapsto0$ 给出的态射。可以算出 $\mathcal F^{(1)}_m$ 是对应于模 $\F_p[T]/T^{p^m}$ 的凝聚层，转移映射由 $T\mapsto T^p$ 诱导（看出这一点的一个快捷方式是：$\varinjlim_mF^{(1)}_m$ 是 $\F_p[[\Z_p]]$ 的庞特里亚金对偶，而 $\varinjlim_m\F_p[T]/T^{p^m}$ 是 $\F_p[[T]]$ 的庞特里亚金对偶，但我们已固定了等同 $\F_p[[\Z_p]]\cong\F_p[[T]]$）。

用 $\F_p[X]$ 的消解
\[
\F_p[[T]][X]\xrightarrow{\ \cdot T\ }\F_p[[T]][X]
\]
计算导出拉回 $\bar\pi_1^*$，可得 $\bar\pi_1^*\mathcal F^{(1)}_m$ 是两项复形，其上同调在次数 $-1$ 为 $\F_p[X]=T^{p^m-1}(\F_p[T]/T^{p^m})[X]$，在次数 0 为 $\F_p[X]=(\F_p[T]/T)[X]$。然而转移映射 $T\mapsto T^p$ 在次数 $-1$ 诱导同构，在次数 0 诱导零映射。由此 $\varinjlim_m\bar\pi_1^*\mathcal F^{(1)}_m$ 是集中在次数 $-1$ 的 $\F_p[X]$，如所断言。
\end{proof}

回到函子 \eqref{eq:a4-7-1-4}，即 $F^\times$（而非 $\Z_p$）的局部解析表示的情形。此时得到函子
\begin{equation}\tag{7.1.10}\label{eq:a4-7-1-10}
\mathfrak A^{\mathrm{rig}}_{F^\times}:D^b_{\mathrm{f.p.}}(\mathrm{an}\,F^\times)\lra D^b_{\mathrm{coh}}(\mathcal X_1).
\end{equation}
从固态局部解析表示的观点看，该函子在 [RJRC25, 4.4] 中有更详细的讨论。

\begin{rem}{7.1.11}\label{rem:a4-7-1-11}
但要注意，函子 \eqref{eq:a4-7-1-10} 在范畴 $D^b_{\mathrm{f.p.}}(\mathrm{an}\,F^\times)$ 上不可能全忠实：粗略地说，该函子涉及从代数 $\mathbb G_m$ 到解析 $\mathbb G_m$ 的解析化。举一个显式的例子：取 $\pi$ 为万有非分歧表示
\[
F^\times\lra F^\times/\mathcal O_F^\times\cong\Z\lra
\Gamma(\mathbb G_m^{\mathrm{alg}},\mathcal O^\times_{\mathbb G_m^{\mathrm{alg}}})=L[T,T^{-1}]^\times,
\]
把 $1$ 映到 $T$，其中 $\mathbb G_m^{\mathrm{alg}}$ 表示概形 $\mathbb G_m$。则 \eqref{eq:a4-7-1-10} 把 $\pi$ 映到刚性解析空间 $\mathbb G_m=\{1\}\times\mathbb G_m\subseteq\mathcal W\times\mathbb G_m=\mathcal T$（其中 $1\in\mathcal W$ 表示 $\mathcal O_F^\times$ 的平凡特征标）在平凡 $\mathbb G_m$-作用下的商上的结构层。由此 $\End_{F^\times}(\pi,\pi)=L[T^{\pm1}]$，而 $\mathcal O_{\mathcal X_1}\widehat\otimes^{L}_{D(F^\times)}\pi$ 的自同态环是刚性解析 $\mathbb G_m$ 上的函数环。事实上，函子 \eqref{eq:a4-7-1-10} 可视为表示 $\pi$ 在 $\mathcal X_1$ 上的某种局部化。
\end{rem}

% 本片止于印刷页 274 页末（§7.1），续片自 PDF 288 页起
